
\chapter{形式语言与自动机}


\section{形式语言}


形式语言(formal language)是用严格符号系统来描述字符串集合的一种数学工具.
在自然语言处理、编译原理和计算理论中, 它常用于刻画句子的结构规则与可计算边界.

通常先给定一个字母表$\Sigma$, 即非空有限符号集合.
由字母表中的符号按顺序连接得到的有限序列称为字符串(string).

\begin{itemize}
\item 空串记作$\epsilon$.
\item 字符串$x$的长度记作$|x|$.
\item 字母表$\Sigma$上的所有字符串构成的集合记作$\Sigma^*$.
\item 除空串外的所有字符串集合记作$\Sigma^+$.
\end{itemize}

如果一个集合$L$满足$L \subseteq \Sigma^*$, 则称$L$是字母表$\Sigma$上的一个形式语言.

例如, 若$\Sigma = \{a,b\}$, 则

\begin{itemize}
\item $\epsilon \in \Sigma^*$.
\item \texttt{ab}, \texttt{aab}, \texttt{bba} 都属于 $\Sigma^*$.
\item $L = \{a^n b^n | n \ge 1\}$是$\Sigma$上的一个语言.
\end{itemize}

\subsection{形式语法}


\textbf{理解形式语法的核心是从非终结符到终结符的规则的集合}

形式语法(formal grammar)通常记为四元组:
\[
G = (N, \Sigma, P, S)
\]
其中:

\begin{itemize}
\item $N$是非终结符集合(nonterminal set).
\item $\Sigma$是终结符集合(terminal set).
\item $P$是一组重写规则的有限集合.
\item $S$是开始符号, 且$S \in N$.
\end{itemize}

非终结符表示可继续展开的语法成分, 终结符表示最终生成的符号串中的实际符号.
产生式给出了从一个符号串到另一个符号串的替换规则.

例如文法
\[
G = (\{S\}, \{a,b\}, \{S \to a S b, S \to a b\}, S)
\]
可以生成语言
\[
L(G) = \{a^n b^n | n \ge 1\}
\]

\subsection{推导}


若文法中存在产生式, 可将符号串$\alpha$中的某部分替换为另一符号串$\beta$, 就称发生了一步推导, 记作:
\[
\alpha \Rightarrow \beta
\]

经过有限步推导, 若能由开始符号得到终结符串$w$, 即
\[
S \Rightarrow^{*} w
\]
则称$w$可由文法$G$生成.

其中:

\begin{itemize}
\item $\Rightarrow$表示一步推导.
\item $\Rightarrow^{*}$表示经过零步或多步推导.
\item $\Rightarrow^{+}$表示经过一步或多步推导.
\end{itemize}

在上下文无关文法中, 常讨论左推导和右推导:

\begin{itemize}
\item 左推导: 每一步总是展开最左边的非终结符.
\item 右推导: 每一步总是展开最右边的非终结符.
\end{itemize}

例如对于文法$S \to a S b | a b$, 字符串\texttt{aabb}的一种左推导为:
\[
S \Rightarrow a S b \Rightarrow a a b b
\]

\subsection{文法类型}


按照乔姆斯基层级(Chomsky hierarchy), 文法可分为四类.

\subsubsection{正则文法}


正则文法(type-3 grammar)的所有产生式通常限制为如下形式:

\begin{itemize}
\item 右线性文法: $A \to a B$或$A \to a$.
\item 左线性文法: $A \to B a$或$A \to a$.
\end{itemize}

其中$A,B \in V_N$, $a \in V_T$.

正则文法生成的语言称为正则语言(regular language).
这类语言结构最简单, 能被有限自动机识别.

典型例子:
\[
S \to a S | b S | \epsilon
\]
它生成所有由$a,b$构成的字符串, 即$\{a,b\}^*$.

\subsubsection{上下文无关文法}


上下文无关文法(context-free grammar, CFG)的产生式具有形式:
\[
A \to \gamma
\]
其中$A \in N$, $\gamma \in (N \cup \Sigma)^*$.

也就是说, 产生式左侧必须是单个非终结符, 因而替换不依赖其上下文.

CFG能够描述嵌套和层次结构, 如:

\begin{itemize}
\item 括号匹配.
\item 算术表达式.
\item 句法树结构.
\end{itemize}

例如
\[
S \to a S b | \epsilon
\]
生成语言$L = \{a^n b^n | n \ge 0\}$, 这是典型的上下文无关语言, 但不是正则语言.

\subsubsection{上下文有关文法}


上下文有关文法(context-sensitive grammar)的产生式一般写作:
\[
\alpha A \beta \to \alpha \gamma \beta
\]
其中$\gamma \ne \epsilon$.

这表示非终结符$A$可以在上下文$\alpha$和$\beta$中被替换为$\gamma$, 替换依赖周围环境, 因此称为上下文有关.

这类文法满足长度不减性质:
\[
|\alpha A \beta| \le |\alpha \gamma \beta|
\]

它能描述比CFG更复杂的约束, 例如多个成分之间的一致性关系.

\subsubsection{无约束文法}


无约束文法(unrestricted grammar)或短语结构文法(type-0 grammar)对产生式几乎没有限制, 一般只要求:
\[
\alpha \to \beta
\]
其中$\alpha$至少包含一个非终结符.

它所生成的语言就是递归可枚举语言(recursively enumerable language), 对应图灵机可识别的语言类.

从表达能力上看, 四类文法满足如下包含关系:

\begin{itemize}
\item type-3 文法包含于 type-2 文法.
\item type-2 文法包含于 type-1 文法.
\item type-1 文法包含于 type-0 文法.
\end{itemize}

\subsection{CFG派生树}


上下文无关文法可以自然地表示成派生树(parse tree)或语法树.

派生树的基本特点是:

\begin{itemize}
\item 根节点是开始符号$S$.
\item 内部节点是非终结符.
\item 叶节点是终结符或$\epsilon$.
\item 从左到右读取叶节点, 得到该树对应的句子.
\end{itemize}

例如对于文法
\[
E \to E + E | E \cdot E | (E) | \operatorname{id}
\]
字符串\texttt{id + id * id}可以对应不同的派生树.
若一个字符串在同一文法下有两棵不同的派生树, 则称该文法是二义的(ambiguous).

二义性意味着同一个句子可能有不同结构解释.
在自然语言处理中, 句法歧义与文法二义性密切相关.

为了消除二义性, 常通过如下方法改写文法:

\begin{itemize}
\item 将不同优先级运算分层定义.
\item 显式区分结合方向.
\item 引入更多非终结符表示中间层级.
\end{itemize}

下面给出一个具体的CFG派生树例子. 对于文法
\[
S \to a S b | a b
\]
字符串\texttt{aabb}可以按如下方式生成:
\[
S \Rightarrow a S b \Rightarrow a a b b
\]

其对应的派生树如下:

\begin{quote}
\small
\textbf{文法$S \to a S b | a b$生成字符串\texttt{aabb}的派生树}
\par 节点: root: $S$; a1: $a$; s1: $S$; b2: $b$; a2: $a$; b1: $b$
\par 边: root -\textgreater{} a1; root -\textgreater{} s1; root -\textgreater{} b2; s1 -\textgreater{} a2; s1 -\textgreater{} b1
\end{quote}

从左到右读取叶节点, 依次得到\texttt{a a b b}, 即字符串\texttt{aabb}.
这个例子说明, 派生树不仅展示了最终生成结果, 也明确展示了中间非终结符是如何逐层展开的.

\section{自动机}


\textbf{自动机是状态机}

自动机(automaton)是对计算过程的抽象模型.
它从某一初始状态出发, 依次读取输入符号, 按照转移规则改变状态, 最终判断输入串是否被接受.

不同能力的自动机对应不同能力的语言类, 这是形式语言理论的核心内容之一.

\subsection{有限自动机}


有限自动机(finite automaton)只有有限个状态, \textbf{常用于识别正则语言}.

确定有限自动机(deterministic finite automaton, DFA)可表示为五元组:
\[
M = (Q, \Sigma, \delta, q_0, F)
\]
其中:

\begin{itemize}
\item $Q$是有限状态集合.
\item $\Sigma$是输入字母表.
\item $\delta: Q \times \Sigma \to Q$是状态转移函数.
\item $q_0 \in Q$是初始状态.
\item $F \subseteq Q$是接受状态集合.
\end{itemize}

自动机从$q_0$开始, 每读取一个输入符号就按$\delta$转移一次.
若输入读完后停在某个接受状态, 则该串被接受.

例如, 构造一个识别“含偶数个$a$”的语言的DFA:

\begin{itemize}
\item 状态$q_e$表示目前读到偶数个$a$.
\item 状态$q_o$表示目前读到奇数个$a$.
\item 读到$b$时状态不变.
\item 读到$a$时在$q_e$与$q_o$之间切换.
\end{itemize}

该自动机识别的语言是:
由所有“包含偶数个\texttt{a}”的字符串组成的集合.

有限自动机的特点是没有额外存储能力, 因此不能识别需要无界计数的语言, 如$\{a^n b^n | n \ge 1\}$.

\subsection{不确定的有限自动机}


不确定有限自动机(nondeterministic finite automaton, NFA)与DFA的区别在于:

\begin{itemize}
\item 对同一状态和同一输入符号, 可以有多个后继状态.
\item 允许存在$\epsilon$转移, 即不消耗输入符号而改变状态.
\end{itemize}

NFA的转移函数通常写为:
\[
\delta: Q \times (\Sigma \cup \{\epsilon\}) \to 2^Q
\]

NFA接受一个串, 是指存在至少一条状态路径使得整个输入串被读完且最终进入接受状态.

虽然NFA看起来比DFA更灵活, 但二者在识别能力上等价:

\begin{itemize}
\item 每个DFA都可以看作特殊的NFA.
\item 每个NFA都可以通过子集构造法(subset construction)转化为等价的DFA.
\end{itemize}

因此, DFA可识别的语言、NFA可识别的语言与正则语言是同一个语言类.

正则表达式、正则文法、DFA和NFA在表达能力上完全等价, 只是描述形式不同.

\subsection{自动机与文法类型对应}


形式语言理论中, 文法类型与自动机模型之间存在经典对应关系:
\begin{table}[H]
\centering
\begingroup
\setlength{\tabcolsep}{2pt}
\small
\begin{tabular}{|p{0.43\textwidth}|p{0.43\textwidth}|}
\hline
type-3 & 有限自动机 \\ \hline
type-2 & 下推自动机(pushdown automaton, PDA) \\ \hline
type-1 & 线性有界自动机(linear bounded automaton, LBA) \\ \hline
type-0 & 图灵机(Turing machine) \\ \hline
\end{tabular}
\endgroup
\end{table}
这种对应关系说明:

\begin{itemize}
\item 文法描述“如何生成”字符串.
\item 自动机描述“如何识别”字符串.
\item 两者从生成和识别两个角度刻画了同一类语言.
\end{itemize}

从自然语言处理角度看:

\begin{itemize}
\item 词法分析通常可由正则语言与有限自动机完成.
\item 句法分析更多依赖上下文无关文法及其扩展.
\item 更复杂的自然语言现象往往超出简单CFG, 需要更强的形式系统来描述.
\end{itemize}

因此, 形式语言与自动机为理解语言结构、设计分析算法以及研究模型能力边界提供了统一的理论基础.
